\section{Introduction}

Firms that lay off the most workers tend to cut wages the least. Why do some of the firms layoff workers instead of cutting wages? In most of the existing literature the answer is ``firms are unable to cut wages for exogenous reasons and thus have to resort to layoffs''. The assumption of exogenous wage rigidity, potentially justifiable via minimum wage, morale costs to wage cuts, or sectoral bargaining, is highly prevalent in macroeconomic literature. It is used to amplify employment outcomes as well as firm responses to monetary shocks. However, recent survey evidence (\textcite{davis2025},\textcite{bertheau2025}) seems to suggest that firms are more free in their choice than previously thought and in fact sometimes actively prefer layoffs instead of potentially available wage cuts. \\
In this paper, I propose an alternative explanation as to why firms layoff workers instead of cutting their wages. Rather than laying off workers because they have to, firms \textit{choose} to layoff workers and then prefer to not cut wages of the survivors. For the core intuition, I refer to the question asked in a firm survey conducted by \textcite{bertheau2025}: ``Why didn't you lower pay instead of laying off employees?". The modal answer -- ``Layoffs give better control over who leaves the firm''. I interpret this answer as firms finding some workers more productive than others and preferring layoffs as a tool to get rid of those workers (as opposed to wage cuts that could incentivize workers to look for a different job). With this in mind, I build a dynamic model with a frictional labor market where workers, both unemployed and employed, look for jobs. Firms may hire workers each period, with each worker-firm match having its own quality, unknown to the parties prior to working. In order to facilitate the idea that layoffs are a better tool of getting rid of poor matches, I assume this quality is only observed by the firm and that the firm cannot wage discriminate within any particular cohort (workers hired in the same period). Layoffs are thus the only tool to affect the quality within a cohort. When firms resort to it, the surviving cohort is now of higher average quality than before, and thus the firm is less incentivized to cut their wages. \\
Beyond the survey evidence, I use the French matched employer-employee data between 2009 and 2019 to discipline and validate my model. First, I confirm the literature finding (\textcite{enrlich2024}) that firms which layoff the most exhibit the most rigid wages. I document the layoff rate and the passthrough of firm-level productivity shocks to the wages of job stayers. I group firms based on their average layoff rates and find that, after a positive productivity shock normalized to $100\%$, wages in firms firing the least increase by $1.7\%$, while wages of firms firing the most change wages by less than $1\%$.
Furthermore, I show that the connection between layoffs and rigid wages exists not only across firms, but also within firms, across worker tenure. Workers with tenure of less than 2 years exhibit a layoff rate of $5-11\%$ as compared to the $2\%$ layoff rate of workers with 4-5 years of tenure. On the other hand, wages of workers with less than 1 year of tenure essentially do not respond to negative productivity shocks, as compared to more senior workers whose wages fall by $2.4\%$ in response to a $100\%$ productivity shock.
%On the other hand, wages of workers with less than 3 years of tenure respond negatively to positive productivity shocks, as opposed to more senior workers. Looking deeper into this, I find that the negative response to productivity shocks come mainly from the negative side: wages of juniors workers rise no matter what, but even faster in response to negative shocks. 
Lastly, I look at relative wage changes between junior and senior workers when a firm begins layoffs. I find that, whenever firms layoff their workers, the wage gap between senior and junior workers contracts. \\
%\textcolor{red}{Past that I start writing the details + related literature} 
%Model description paragraph
To understand these facts, this paper builds an equilibrium model of the labor market with search frictions and one-sided limited commitment on the worker side. Firms employ a continuum of workers, exhibit decreasing returns to scale, and are subject to firm-level productivity shocks. Search is directed, firms post dynamic contracts to attract workers. The contracts comprise of future productivity history-contingent wages and layoff risk. Workers are risk-averse and cannot commit to stay at their job. On the other hand, firms are able to credibly promise to insure their workers against future productivity shocks. Firms thus face a tradeoff between insuring workers against future losses and incentivizing workers to stay (leave) during times where they're the most (least) productive. Lastly all the worker-firm matches, while ex-ante homogenous, are subject to an initial and permanent shock upon beginning of employment. While the firm-level productivity shock is common knowledge, the match-specific shock is observed only by the firm. \\
%1 paragraph explaining how I make this model tractable. Do I talk about the main assumption before or after? I think after. But I mention quality before
Firms exhibit decreasing returns to scale and thus handle the contracts for each employed worker together. In principle, this makes the model untractable: in recursive formulation, the firm would need to track the contracts of the continuum of workers, resulting in an infinite-dimensional state-space. Luckily, due to the directed nature of the search, all the workers hired by the firm in a single period will be hired on the same contract. It is thus sufficient to keep track of contracts just for each cohort that the firm employs. This discretizes the state-space and, for firms with finite age, makes it similarly finite: firm of age 10 will employ at most 10 different cohorts. Furthermore, up to an approximation, the state-space can be bounded: I show both theoretically and empirically that wages in workers' contracts tend to converge. I use this to justify putting all the workers past the certain tenure level on the same contract as those contracts are already quite similar. Thus, instead of optimizing a continuum of contracts, it is enough to deal with a finite number of tenure-specific contracts. \\
I assume that firms are unable to wage discriminate based on quality within a cohort. This leaves layoffs as the only way to affect quality within a cohort. During low productivity states, layoffs shoot two birds with one stone: they get to rid of overpaid (relative to the productivity) labor and cleanse the workforce of poor matches. In whichever cohorts that suffered layoffs, the surviving workers are of higher average quality. The firm is now incentivized to keep the survivors and thus will not cut their wages by as much as the other cohorts, where the proportion of good matches did not change. This mechanism is at the key of this paper's view on layoffs and pay cuts: whenever and wherever the firm chooses to fire workers, it will not want to cut wages of the survivors. \\
Unlike models of exogenous wage rigidity, this paper can explain the connection between layoffs and wage cuts not just between, but within firms. When the firm chooses who to layoff, it is optimal to get rid of poor matches who are the cheapest to fire. In my model, this is almost always junior workers. Therefore, the model exhibits relatively larger layoff rates for the most junior workers in each firm. The individual layoff risk depends not only the worker's absolute tenure, both on their relative tenure within the firm: the cheapest, and thus most junior workers get cut first. This is consistent with findings that workers who are hired last are likely to be fired first, like \textcite{buhai2014}. After layoffs, the firm is not as incentivized to cut wages for surviving juniors as for seniors whose average quality has not improved. The model can thus replicate the heterogeneity in layoffs and wage rigidity across tenure documented in the data.
\\
The assumption of no match quality-based discrimination is at the core of my model. The intuition that the firms provide: ``Layoffs give better control over who leaves the firm" -- requires that firms consider layoffs somehow more useful than wage cuts when it comes to getting rid undesirable matches. The same model with complete information about match quality would have no special role for layoffs and would be rendered useless for generating meaningful layoffs. I take several approaches to justify this modeling choice. One interpretation of this assumption is nondiscrimination laws: if firms are unable to justify why they `prefer' one worker other another, they cannot be allowed to selectively change wages. This is different from commonly observed differences in worker quality: there is no doubt in productivity differences when Ann publishes significantly more papers than Bob, and I do not need to restrict firms from offering contracts specific to observable quality. Indeed, it is enough to assume existence of some preference that the firms holds over its workers that the workers are unaware of. \\ 
Within the case of asymetrically observed differences in match quality, I consider an extension where the firm is allowed to wage discriminate however it sees fit. There is still a reason not to do that though: as long as workers do not know their own quality, they're essentially ``risk-sharing'' whatever layoff risk comes their way. This makes layoffs cheaper for ther firm, and the firm thus has two separate and opposing tools for getting rid of poor matches. In Appendix~\ref{microfoundation} I consider a signalling game, where firm's wage choices affect workers' beliefs about their own quality. I find that, under sufficiently inelastic probability of retaining workers from finding jobs elsewhere, the pooling equilibrium, where the firm does not reveal the matches types, exists. I can then justify my assumption as focusing on the pooling equilibrium of the extended model.

%Quantitative paragraph(s)

%Assumptions discussion paragraph? For example, if I get the results for the pooling equilibrium, I could maybe write a paragraph on that

\textbf{Related Literature} \\
To be done